\subsection{감사 AUD3: VKGAP 절과 NUM 절}\label{AUD3:sec}

\paragraph{범위와 방법.}
감사 대상은 \texttt{agents/VKGAP/frag.tex}(비공명 $\Rightarrow$ 약수 간격, VKGAP-A/B/C)와 \texttt{agents/NUM/frag.tex}(다중집합 BFS와 공명 높이 실험)이다.
편집 전 원본은 각각 \texttt{frag.tex.orig\_before\_AUD3}로 보존하였다. NUM의 \texttt{frag.tex}는 원래 \texttt{frag\_src.tex}와
\texttt{code/assemble\_frag.py}로 생성되지만, 지시에 따라 \texttt{frag.tex}를 직접 고쳤다(\texttt{frag\_src.tex}는 손대지 않음; 재생성하면 AUD3 수정이 사라진다).
모든 편집은 \texttt{\% AUD3:} 주석(이유 포함)으로 표시하였고, 편집 후 두 파일 모두 \texttt{test\_frag.sh}에서 \texttt{COMPILE OK}이다.
VKGAP의 \PROVED, \REFUTED{} 항목과 모든 조건부 함의를 처음부터 다시 유도하였다(각 부등식과 상수를 손으로, 그리고 아래 스크립트로 다시 계산).
NUM의 알고리즘 정당성 논증을 다시 유도하였고, 주장된 값은 NUM 코드와 독립적으로 새로 작성한 코드로 점검하였다.
스크립트는 \texttt{agents/AUD3/code/}, 출력은 \texttt{agents/AUD3/out/}에 있다.
\begin{itemize}[leftmargin=1.5em]
\item \texttt{aud3\_vkgap\_constants.py} $\to$ \texttt{aud3\_vkgap\_constants.txt}: VKGAP의 모든 명시 상수(분산 하한, $c_2$, $128/c_2$, Mellin 한계, $\int|\tilde w|$,
 $-\zeta'(2)/\zeta(2)$, 핵심 명제의 다섯 오차항, $0.375$, $1.555$, $0.2033$, 누적량 나머지 $97/288$의 무작위 시험, (C2), Dirichlet 부등식 $4K^M\le(8/\eta)^M$, 조화합 한계).
\item \texttt{aud3\_vkgap\_primes.py} $\to$ \texttt{aud3\_vkgap\_primes.txt}: $L=10^4,10^5,10^6$에서 $A_L(\tau)$, $\sum_P(1-\cos)$ (Lipschitz 인증서 포함), 쌍 논법 비, 매끄러운 합;
 공약수 보조정리의 정수 전수 검사; 작은 $L$의 약수 간격.
\item \texttt{aud3\_num\_dp.py} $\to$ \texttt{aud3\_num\_dp.txt}: $L\le18$에서 서로 다른 약수 0/1 DP와 다중집합 BFS(둘 다 새 코드)의 점별 비교, $H$, 히스토그램, $a_c$, 이동 항등식; 최적 표현 검증.
\item \texttt{aud3\_mitm.c} $\to$ \texttt{aud3\_mitm\_L23.txt}, \texttt{aud3\_mitm\_L25.txt}, \texttt{aud3\_mitm\_L25\_a8.txt}: $S_3$ 비트집합에 의한 소속 판정(새 코드).
\item \texttt{aud3\_num\_res.py} $\to$ \texttt{aud3\_num\_res.txt}: 공명 높이 21칸 재계산, $K/M$, LLL 점 22개 재평가, $\tau_0(\eta)$; \texttt{aud3\_ratefn.py} $\to$ \texttt{aud3\_ratefn.txt}: $I(x)$, Dirichlet $Q$.
\end{itemize}
모든 작업은 \texttt{nice -n 10}, 메모리 3\,GB 이하(최대: $L=25$ 검사의 $2.84$\,GB 비트집합)로 수행하였고, 최종 실행은 모두 5분 이내이다(첫 소수합 시도 하나는 약 9분 뒤 중단하고 범위를 줄여 다시 실행). $L=23,25$ BFS는 다시 돌리지 않았다.

\paragraph{총평.}
\textbf{VKGAP}: 29개 항목 중 CONFIRMED 24, MINOR-FIX 5, GAP 0, WRONG 0.
\textbf{NUM}: 15개 항목 중 CONFIRMED 14, MINOR-FIX 1, GAP 0, WRONG 0.
VKGAP의 수학적 핵심(공약수 보조정리, 매끄러운 VK 소수합, 격자 공명 구간의 적분 처리, Fejér 국소 극한과 명시 상수, 전 구간 덮기, greedy 장부, 일반 전이, RH 판, Dirichlet 천장)은
적대적으로 다시 유도해도 그대로 성립한다. 특히 정리 VKGAP-A (\texttt{VKGAP:thm-A})와 VKGAP-B (\texttt{VKGAP:thm-B})는 교과서 입력 (E1)--(E4)에 조건부인 완전한 증명으로 확인된다.
수정은 모두 수치 오기, 반박 문턱, 태그 표기, 정의역 세부, 기억된 공식에 대한 단서에 관한 것이다.
NUM의 알고리즘 정당성(중복 제거 $\Rightarrow$ 다중집합 BFS)과 계수 보조정리는 옳고, 독립 코드로 점검한 모든 수치가 일치한다.
AUD3는 추가로 $\kappa_{25}(a_8)\ge8$ (즉 $H(25)\ge8$)을 BFS와 독립으로 확인하였다(명제~\ref{AUD3:prop-a8}).

%----------------------------------------------------------------------
\subsubsection{VKGAP 판정표}\label{AUD3:ssec-vk-table}
\begin{center}\small
\begin{longtable}{p{0.22\textwidth}p{0.11\textwidth}p{0.12\textwidth}p{0.47\textwidth}}
\toprule
항목 (VKGAP 라벨) & 원 태그 & 판정 & 이유\\\midrule
\endhead
Fejér 항등식, $F_T\ge4T/\pi^2$ ($|x|\le\frac1{2T}$) & PROVED & CONFIRMED & 직접 적분 $\frac{1-\cos2\pi Tx}{2\pi^2Tx^2}=F_T$; $\sin v/v\ge2/\pi$.\\
\texttt{lem-cofactor} & PROVED & CONFIRMED & 홀수 $e$: $2^b<p^a<2^{b+1}$, $e'=e2^b/p^a\in(e/2,e)$, $2^b<p^a\le L\Rightarrow b\le a_2=v_2(Q_y)$ ($y\ge16\Rightarrow2\notin P_y$). $P_y$의 소수가 $Q_y$에 남는 경우($y\le\sqrt L$)도 논증은 $p^a\mid\Lam_L$만 쓴다. 수치: 정수 전수 검사 156개 $(L,y)$에서 위반 0.\\
\texttt{lem-phi1} & PROVED & CONFIRMED & $(1-\cos x)/x^2$가 $(0,\pi]$에서 감소.\\
\texttt{lem-pair} & PROVED (+E1) & CONFIRMED & $\sum_{p,p'}(x_p-x_{p'})^2=2M^2\mathrm{Var}$, $|\tau(\log p-\log p')|<\pi$ ($|\tau|\le\pi/\log2=4.532$), $0.19^2(\log1.5)^2=0.0059349$. 수치 비 $0.0173\ge\frac2{\pi^2}\mathrm{Var}=0.0080$.\\
\texttt{lem-mellin} & PROVED & CONFIRMED & 부분적분식 재계산(직접 적분과 $10^{-17}$ 이내 일치), $4/(\kappa|s||s+1|)\le80/t^2$, $4\sqrt{80}=35.78$; $\int|w-1_{[1/2,1]}|=\kappa$. 격자 검사 최대비 $0.74\le1$.\\
\texttt{lem-VK} & COND (E1)--(E4) & CONFIRMED & 아래 메모 (V1).\\
\texttt{cor-NRuncond} & COND & CONFIRMED & $0.45-\frac{1.5}{\sqrt{17}}-0.05=0.03620$; $0.035/0.51=0.0686\ge1/15$. \emph{하한만} 필요하다는 주장 옳음 ($1-\cos\ge0$, $1_{P_y}(p)\ge w(p/y)\log p/\log y$). 수치: $L=10^4,10^5$에서 $[4,2000]$, $[4,600]$ 인증 하한 $\sum(1-\cos)/M\ge0.29$.\\
\texttt{rem-F}(b) & PROVED & CONFIRMED & $5-4\cos x\le1+2x^2$; $\sqrt{(1+a)/2}=0.990178$; $(1-a)/2=0.019547$.\\
\texttt{rem-F}(c) & NUMERIC & CONFIRMED & $|F(1)|=0.9805724$; 격자 최대 $=|F(1)|$; 단조성 인증 논법 점검.\\
\texttt{lem-cumulant} & PROVED & CONFIRMED & $\frac16+\frac{49}{288}=\frac{97}{288}$, $|z|\le\frac7{48}$, $|\log(1+z)-z|\le|z|^2$; 무작위 시험 20000회 최대비 $0.49$.\\
조건 (C1)--(C3) & PROVED & CONFIRMED & $n$에 대한 단조성($n\ge2(\log y)^2/c_2$, $n\ge1/\eta'$), $n=n_0$에서 $8\log\log y\ge6.91+\frac12\log K+2\log\log y$.\\
\texttt{prop-core} & PROVED & MINOR-FIX & 모든 상수 재계산 일치(메모 (V3)). 다만 (d)가 쓰는 \texttt{lem-pair}의 $v_y\ge v_0$는 (E1)에 의존하므로 원장 태그를 ``\PROVED{} (+(E1))''로 맞춤(본문 진술은 이미 이를 밝힘).\\
\texttt{rem-core-quant} & PROVED & CONFIRMED & $F_T\le T$에서 $\Pr\ge0.17/(T\sqrt V)$.\\
\texttt{thm-A-y} & COND & CONFIRMED & 메모 (V4).\\
\texttt{thm-A} (VKGAP-A) & COND (E1)--(E4) & CONFIRMED & $2r_L=\delta$, $e^{-\delta}\ge1-\delta$, $e^{\delta}\le1+2\delta$; $X_0=e^{\sigma(L)}$; $128/c_2=107060<1.08\cdot10^5$.\\
\texttt{thm-B} (VKGAP-B), $N(b)$ & COND & CONFIRMED & 메모 (V5).\\
\texttt{lem-greedy} & PROVED & CONFIRMED & $1-e^{-x}\le x$로 $R'\le Re^{-G}$; 블록당 $\le2^j/G_j+1$; 위쪽은 $\ell^*$ 증가; 끝단 $\le a/\log2+1$, (K3)로 $\le k+2$; 서로 다름은 (K2).\\
\texttt{cor-regimes} & PROVED & MINOR-FIX & (a),(b),(c) 계산 옳음 ($\sum_{k\le K}\frac1{1+k\log2}\le1+\frac{\log(1+\log L)}{\log2}$ 수치 확인). (c)의 $G$는 $D>L$에서 감소하므로 보조정리의 ``비감소'' 가정을 위해 $G(D):=G(L)$ ($D\ge L$)로 둔다는 단서 추가(결과 불변).\\
\texttt{lem-smallzone} & COND & CONFIRMED & $\sigma(y(\ell))=\ell$, $(\log y)^3\ge\frac{c_2}{256}\frac{\ell}{\log\ell}$, $\sigma(L)\le\ell\le\ell_1$에서의 비교, $2\sum_j2^{j/2}\le6.83\sqrt{\ell_1}$; 아래쪽 끝 $R<e^{\ell_{\min}}\le L$은 그 자체가 약수(큰 $L$).\\
\texttt{def-NR} (NR$\Rightarrow$NR$'$) & PROVED & CONFIRMED & $M-\mathrm{Re}A\ge M-|A|$.\\
\texttt{thm-transfer} & COND & CONFIRMED & 메모 (V6).\\
\texttt{cor-three}(1) & COND & CONFIRMED & $D_*<\ell_1$이므로 실제 하위항은 $O((\log L)^{3/2}(\log\log L)^4)$ (주~\texttt{rem-B-honest}와 일치).\\
\texttt{lem-RH}, \texttt{cor-three}(2) & COND RH & MINOR-FIX & 유도 옳음(메모 (V7)). 태그에 (E1),(E3),(E4) 추가; 주~\texttt{rem-RH-check}에 lead가 기억한 절단 공식의 오차항이 틀렸다는 단서 추가(증명에는 쓰이지 않음).\\
\texttt{cor-three}(3) (VKGAP-C) & COND NR & CONFIRMED & $\psi(L)/(cM-4)=O(\log L/c)$; 약한 $\mathrm{NR}'(e^{cL/(\log L)^2},\eta)$로 충분; $c_4=\min(c/3,\eta/16)$.\\
\texttt{rem-cutoff} & COND (E1) / REFUTED & MINOR-FIX & 점근식 옳음. 반박 문턱은 $\eta>1-|F(1)|=0.0194276$이어야 한다: $0.0194<\eta\le0.019427$이면 $1-\eta\ge|F(1)|$라 $\tau=1$ 논법이 반박을 주지 않는다. 본문과 원장 수정.\\
\texttt{prop-dirichlet} & PROVED & CONFIRMED & 비둘기집, $|e^{2\pi i\theta}-1|\le2\pi\|\theta\|$, $2\sin^2(\pi\|\theta\|)\le2\pi^2\|\theta\|^2$; $K\le(2\pi+1)/\eta$이고 $(8/7.283)^{15}\ge4$. ($\eta$ 격자 수치로는 $M\ge13$에서 이미 성립.)\\
\texttt{rem-lattice} & HEUR/NUMERIC & MINOR-FIX & 극한 분산 $2-\log^22-2\log2-(1-\log2)^2=0.039094$ (원문 $0.03907$은 오기; NUM 명제 \texttt{prop-var}와 일치하도록 수정). (N1)의 같은 값도 수정.\\
(N1)--(N3) & NUMERIC & CONFIRMED & 독립 재계산: $\mathrm{Var}_P=0.03937/0.03932/0.03916$, 쌍 비 $0.0174/0.0173/0.0173$, $\max|A|/M=0.714/0.715/0.716$ ($\tau=4$), 격자 공명 $0.00914/0.00585$, $1/G_{1/2}$ 값 일치. 단 (N1)의 $\min\sum(1-\cos)/M$는 성긴 격자값이며 간격 $0.005$에서는 $0.316$ ($L=10^4$)로 조금 더 작다(주장된 용도에는 무관).\\
\texttt{rem-ladder}, \texttt{rem-fixedsize}, \texttt{rem-A-new} & HEUR/OPEN & CONFIRMED & 태그대로(격상 없음). $L^1$-Fourier 천장 논증은 발견적으로 타당.\\
\bottomrule
\end{longtable}
\end{center}

%----------------------------------------------------------------------
\subsubsection{VKGAP 재유도 메모}\label{AUD3:ssec-vk-notes}

\paragraph{(V1) 외부 입력과 \texttt{lem-VK}.}
(E1) $\psi(x)=x+O(xe^{-c\sqrt{\log x}})$, $\pi(x)=x/\log x+O(x/\log^2x)$, $\theta(x)\le x\log4$; (E2) VK 영역 $\sigma\ge1-c_1(\log(|t|+3))^{-2/3}(\log\log(|t|+3))^{-1/3}$
(작은 $|t|$로의 확장: $|\gamma|\le t_0$인 영점은 유한개, $\beta<1$, 실수축 $\zeta(\sigma)<0$); (E3) Titchmarsh 9.6(A)형 부분분수와 $N(t+1)-N(t)\ll\log t$,
그리고 $\{1/2\le\sigma\le2,|t|\le2\}$에서 $\zeta'/\zeta+1/(s-1)$의 유계성(첫 영점 높이 $14.13$); (E4) Mellin 역변환 --- 모두 교과서 진술로서 정확하다 (recalled, \EXT).
윤곽: $|{\rm Im}\,z|\le2|\tau|+y^3\le3T_*(y)$ ($T_*\ge y^3$), ${\rm Re}\,z\ge1-\eta_*/2>1-\eta(H+1)$이므로 영점 없음; 유수 $\tilde w(1-i\tau)y^{1-i\tau}$;
$|z-\rho|\ge\sigma_1-\beta>\eta_*/2$; 상수 $36/2\pi=5.73$, $2\cdot1.05\cdot80/2\pi=26.7$, $0.6\cdot160/2\pi=15.3$, $-\zeta'(2)/\zeta(2)=0.5700$ 재계산.
점근: $(2\lambda)^{2/3}=2^{2/3}\log y(\log\log y)^{-2}$에서 $\eta_*\log y\ge\frac{c_1}2(\log\log y)^{5/3}$, 손실 $Z_*\ll(\log y)^{5/2}=e^{\frac52\log\log y}$;
지수 $5/3>1$이므로 $\eps_3(y)\to0$. (수렴은 매우 느리다: 지수가 음이 되려면 $(\log\log y)^{2/3}>10/c_1\ge100$이어야 한다 --- 정리는 존재 진술이므로 무관.)
소수 거듭제곱 $\le\psi(y)-\theta(y)\le3\sqrt y$. 결론: 진술 그대로 옳다.

\paragraph{(V2) 격자 공명 구간과 문턱.}
$\tau_1=1/(2\log y)<|\tau|\le4$에서 $|\varphi|\le e^{-c_2n\tau^2}$ (\texttt{lem-pair}); Fejér 가중치 $\le1$이므로 기여는
$\frac1\pi\int_{\tau_1}^4e^{-c_2n\tau^2}d\tau\le\frac4\pi e^{-c_2n\tau_1^2}$이고 높이 $T$의 인자는 \emph{붙지 않는다} --- 주장 옳음.
이 항이 $10^{-3}V^{-1/2}$ 이하가 되려면 $\frac{c_2n}{4(\log y)^2}\ge\log(10^3\sqrt n\log y)$, 즉 $n\asymp(\log y)^2\log\log y$ ($n_0=\frac{32}{c_2}(\log y)^2\log\log y$)이면 충분하다.
작은 $\tau$ 구간 $|\tau|\le\tau_1$에서는 $|\tau|b=1/2$이므로 누적량 보조정리가 적용된다.

\paragraph{(V3) 국소 극한 상수.}
$G\ge(2\pi)^{-1/2}e^{-1/(2\cdot10^4)}=0.398922$; 오차 (a) $\frac1{50\pi}=0.006366$, (b) $\frac1{200\pi^2}=0.000507$, (c) $\frac{4}{100\pi}=0.012732$
($\int|\tau|^3e^{-\tau^2V/4}=16/V^2$), (d) $\frac4\pi10^{-3}=0.001273$, (e) $0.002$; 합 $0.022878$, 하한 $0.376044\ge0.375$.
상한 국소 평가: $\int e^{-2\tau^2V/\pi^2}=\pi^{3/2}/\sqrt{2V}=3.93740V^{-1/2}$, 합 $3.95797$, $B\le\frac\pi8\cdot3.95797=1.55429\le1.555$.
꼬리 $\frac{2\cdot1.555}{\pi^2}(\frac{\pi^2}6-1)=0.203224\le0.2033$; 여유 $0.375-0.2033=0.1717>0$. 모두 일치 (\texttt{aud3\_vkgap\_constants.txt}).

\paragraph{(V4) 공약수 이동과 전 구간 덮기.}
\texttt{thm-A-y}(i): $q=\max(q_0,\min(\frac12,s/\theta_y))$이면 $|s-q\theta_y|\le1$이고 $n\ge q_0M_y/2=n_0$ ($q(1-q)$는 $(0,\frac12]$에서 증가);
$e^{n/30}\ge T_*$이므로 $T=T_*/2\pi$, $2/T=r_y$. (ii): $z=u-q_0\theta_y\ge2-r_y>0$ ($q_0\theta_y+2\le\sigma(y)$);
$z\le\log Q_y$이면 $s\in[q_0\theta_y,q_0\theta_y+\log2)$, 아니면 $c=Q_y$이고 $s\le\frac12\theta_y+1$ ($\psi(L)\ge\theta_y+\theta(y/2)\ge\theta_y+2r_y$);
위쪽 절반은 $\ell^*=\log\Lam-\ell$와 $d\mapsto\Lam/d$로 $d_-=\Lam/e_+$, $d_+=\Lam/e_-$; 진약수. $[X_0,\Lam/X_0]$ 전체가 덮인다. $\delta=8\pi/T_*(L)$, $X_0=e^{\sigma(L)}$ 확인.

\paragraph{(V5) VKGAP-B 장부.}
(K4) (p1p4 명제 5.18)의 증명을 다시 확인: 절반 감소 $\le\lceil\log_2X_0\rceil$, 중간 $\le\log\Lam/\log(1/\delta)+1$, 끝 $\le k+2$ (서로 다름은 (K2)).
$\log(1/\delta)=(\log L)^{3/2}(\log\log L)^{-3}-\log8\pi$, $\log X_0=O((\log L)^3\log\log L)=o(L(\log\log L)^3(\log L)^{-3/2})$, $\psi(L)=L(1+O(e^{-c\sqrt{\log L}}))$이므로
$H(L)\le(1+O(\frac{(\log\log L)^3}{(\log L)^{3/2}}))\frac{L(\log\log L)^3}{(\log L)^{3/2}}$. $N(b)\le2H(L_b)$, $L_b\le(1+o(1))\log b$, $f(x)=x(\log\log x)^3(\log x)^{-3/2}$의 증가성으로
$N(b)\le(2+o(1))\log b(\log\log\log b)^3/(\log\log b)^{3/2}$. 옳다. (Hughes 경계보다 약하다는 VKGAP의 정직한 비교도 옳다.)

\paragraph{(V6) 일반 전이와 VKGAP-C.}
$q=\ell/(4\theta_P)\le\frac12$, $n\ge\ell/(8\log L)\ge n_0(L)$ ($\ell\ge\ell_1$), $z=\frac34\ell\mp2/T\le\frac38\psi(L)+3\le\log Q$; $\delta(\ell)=4/T$,
$\log(1/\delta)=\min(\log T'-\log2\pi,\eta n/2)-\log4\ge\min(\log T'-2,\frac{\eta D}{16\log L})-2$ ($D\le\ell$).
(ii): $D\in[\ell_1,D_*]$에서 $G\ge\eta D/(32\log L)$ ($\ell_1\ge64\log L/\eta$), 블록당 $\le32\log L/\eta+1$, 블록 수 $\le\log_2L$ (양쪽) $\Rightarrow\frac{64}{\eta\log2}(\log L)^2$;
$D\ge D_*$에서 직접 셈 $\psi(L)/(\log T'-4)+2$; 가장자리는 \texttt{lem-smallzone}. VKGAP-C: $\log T'=cM\ge cL/(3\log L)$.
Erdős \#18 Q1 귀결은 BRIEF의 $h(\Lam_L)=H(L)$와 $\Lam_L$의 practical성(p1p4 Thm.~5.4)에 기댄다. 옳다.

\paragraph{(V7) RH 판.}
$\sigma_1=\frac12+\frac1{\log y}$, RH로 ${\rm Re}\,z>\frac12$에 영점 없음, $|z-\rho|\ge1/\log y$, $|G|\ll\log y\log(|\tau|+y)$, 왼쪽 변 $\frac{36e}{2\pi}\sqrt y Z_{\rm RH}$, 수평 변 $38.2Z_{\rm RH}y^{-4}$.
$|\tau|\le T_{\rm RH}$이면 $\log(|\tau|+y)\le2c_{\rm RH}\sqrt y/\log y$, 여유 $0.03580/0.51=0.0702\ge1/15$. 결과 $H(L)\le(c_{\rm RH}^{-1}+o(1))\sqrt L\log L$ 옳음.
lead의 기억된 절단 공식 오차 $O(x^{1/2}\log x\log(|t|+2))$는 $t=0$에서 RH 아래 $\psi(x)-x\ll\sqrt x\log x$를 뜻하여 알려진 결과(von Koch의 $\sqrt x\log^2x$)보다 강하다; 올바른 크기는
$\sqrt x\log x\log(x(|t|+2))$ 정도 (recalled). VKGAP의 증명은 이 공식을 쓰지 않으므로 결론에는 영향이 없다.

\paragraph{VKGAP 수정 목록 (파일 내 \texttt{\% AUD3:}).}
(1) \texttt{rem-lattice}와 (N1): 극한 분산 $0.03907\to0.039094$. (2) \texttt{rem-cutoff}와 원장: 반박 문턱 $\eta>0.0194\to\eta>1-|F(1)|=0.019427\ldots$.
(3) 원장: \texttt{prop-core}를 ``\PROVED{} (+(E1))''로. (4) \texttt{cor-regimes}(c): $G(D):=G(L)$ ($D\ge L$) 단서. (5) \texttt{cor-three}(2): 태그 \COND{RH 및 (E1),(E3),(E4)}.
(6) \texttt{rem-RH-check}: 기억된 공식에 대한 단서. VKGAP의 \texttt{summary.md}에도 문턱 ``$\eta>0.0194$''가 남아 있다(감사자는 .md를 쓸 수 없으므로 lead가 고칠 것).

%----------------------------------------------------------------------
\subsubsection{NUM 판정표}\label{AUD3:ssec-num-table}
\begin{center}\small
\begin{longtable}{p{0.22\textwidth}p{0.11\textwidth}p{0.12\textwidth}p{0.47\textwidth}}
\toprule
항목 (NUM 라벨) & 원 태그 & 판정 & 이유\\\midrule
\endhead
\texttt{prop-bfs} & PROVED (dedup) & CONFIRMED & 부분합 $\le a<\Lam$이므로 절단 무손실. 중복 제거 정리를 다시 유도: $2d\le a<\Lam$이면 $d\ne\Lam$이라 최소 결손 소수 $p$가 존재; $p=2$면 $2d\mid\Lam$; $p\ge3$면 $r=\frac{p+1}2$의 소인수는 $<p$이고 $q^{v_q(r)}\le r\le L$이므로 $r\mid d$, $pd/r+d/r=2d$, 둘 다 진약수, $\Phi=\sum d_i^2$가 $(u-v)^2/2\ge2$ 증가 (두 경우 모두 $\Phi$ 증가, $\Phi\le a^2$) $\Rightarrow$ 종료. 옳다.\\
\texttt{lem-ac} (1)--(3) & PROVED & CONFIRMED & $1\in D^*$, $\Lam/2\in D^*$; $a_c=\min([0,\Lam)\setminus R_{c-1})$ 논증 옳음.\\
\texttt{lem-shift} (계수 보조정리) & PROVED & CONFIRMED & $U_k\subseteq V_{k-1}$, 크기 일치 $\Rightarrow$ 모든 $k\ge1$에서 $U_k=V_{k-1}$ $\Rightarrow$ $\kappa(\Lam/2+b)=1+\kappa(b)$ ($m=\kappa(b)\ge1$이면 $k=m,m+1$을, $b=0$이면 $k=1$을 쓰면 된다). 표~\texttt{tab-L23}, \texttt{tab-L25}의 $Z^+_k=Z^-_{k-1}$ 직접 확인. AUD3 DP로 $L\le18$에서 점별 성립 확인.\\
\texttt{prop-valid} & NUMERIC & CONFIRMED & AUD3 새 코드(0/1 DP와 다중집합 BFS)가 $L\le18$ (서로 다른 $\Lam$ 11개)에서 점별 일치하고 $H$, 히스토그램, $a_c$가 \texttt{H\_table\_bfs.csv}와 같다. $L=23$: $|S_3|=80\,181\,647=|R_3|$ 재현.\\
\texttt{thm-H23} & NUMERIC & CONFIRMED & 7항 표현(합, $\Lam$의 진약수, 서로 다름) 검증; $a_7\notin S_3+S_3$ ($\kappa(a_7)\ge7$), $a_6\notin S_3+S_2$, $a_5\notin S_3+S_1$, $a_4\notin S_3$, 각 $a_c-1\in R_{c-1}$ --- 새 코드로 재확인. $H(23)\le7$ 자체는 BFS(두 구현)에만 근거.\\
\texttt{thm-H25} & NUMERIC & CONFIRMED (보강) & 8항, 7항 표현 검증; $a_7\notin S_3+S_3$, 표본 6개 ($\kappa=8$인 수 $-\Lam/2$) $\notin S_3+S_3$, $a_6\notin S_3+S_2$, $|S_3\cap[0,\Lam/2)|=242\,846\,840$ (표의 $\Lam/2-Z^-_3$와 일치). 추가로 $\kappa(a_8)\ge8$을 BFS 독립 확인(명제~\ref{AUD3:prop-a8}); NUM 본문 (4),(8)과 원장에 반영.\\
표 \texttt{tab-H}, \texttt{tab-frac} & NUMERIC & CONFIRMED & $L\le18$ 행 재현; $H/\ln L$, $\ln(\Lam-1)/\ln\tau$ 열 재계산 일치.\\
\texttt{rem-H25}, \texttt{rem-H-trend} & HEUR & CONFIRMED & 수치 ($2.49$, $2.46$, $2.23$, $7/2.963=2.36$) 재계산 일치; 태그 적절.\\
\texttt{prop-t1} & PROVED (EXT PNT) & CONFIRMED & 부분적분과 $1/\log u=1/\log L+O(1/\log^2L)$; $|F(1)|=0.9805724$; 문턱 $1-|F(1)|=0.019428$ 정확(VKGAP의 원래 $0.0194$보다 정확).\\
\texttt{rem-t1} & NUMERIC & CONFIRMED & $\tau_0(0.1,0.2,0.3)=2.2988,3.3082,4.1293$ 재현; $g_L(1)/M$ 범위는 NUM 자료 및 AUD3 값($0.98043$, $0.98046$, $0.98054$)과 일치.\\
\texttt{lem-grid} & PROVED & CONFIRMED & $|e^{ix}-1|\le|x|$, $||x|-|y||\le|x-y|$; $|\ell_p|\le\frac{\ln2}2$. 실측 $K/M=0.17$--$0.21$.\\
\texttt{tab-res}, \texttt{prop-res} & NUMERIC & MINOR-FIX & ``격자'' 열 21칸($L\le101$, 세 $\eta$)을 새 코드로 재계산하여 모두 $10^{-3}$ 상대오차 이내 일치; LLL 점 22개 중 $t\le1.2\cdot10^{14}$인 21개는 60자리 재평가로 $g\ge(1-\eta)M$ 확인. 그러나 CSV는 $t$를 유효숫자 15자리로만 저장하여 $t\gtrsim10^{14}$의 LLL 점은 CSV에서 재현 불가($L=373$, $\eta=0.05$: 재계산 $g/M=0.21$). 상한 자체는 실행 중 90자리 평가로 타당하나 재현성 단서를 표 설명에 추가.\\
\texttt{rem-res-heur} & HEUR & CONFIRMED & $I(0.7,0.8,0.9)=0.5760,0.8153,1.1963$ (기울기 $0.250,0.354,0.520$), $Q=8,10,14$ 재현; 태그 적절.\\
\texttt{prop-F} & NUMERIC & CONFIRMED & $|\frac{d}{dt}|F|^2|\le\frac{4\ln2}{1+t^2}+\frac{18t}{(1+t^2)^2}\le5.89<6$ ($t\ge1$), $3/\sqrt{1+2.87^2}=0.98710$; VKGAP의 인증 수치와 일치.\\
\texttt{prop-var} (1),(2) & PROVED (EXT) / NUMERIC & CONFIRMED & 밀도 $2e^{-u}$, 평균 $1-\ln2$, 이차 적률 $2-2\ln2-\ln^22$, 분산 $0.0390940$; 콤팩트 지지에서 약수렴 $\Rightarrow$ 적률 수렴. 쌍 부등식은 사실 증명된 부등식이며 수치 최소비 $2.42$는 $t\to0$ 극한 $\pi^2/4=2.47$과 부합.\\
\bottomrule
\end{longtable}
\end{center}

%----------------------------------------------------------------------
\subsubsection{NUM 재검증 메모}\label{AUD3:ssec-num-notes}

\begin{proposition}[$H(25)\ge8$의 BFS 독립 확인]\label{AUD3:prop-a8}
\NUMERIC\ (\texttt{agents/AUD3/code/aud3\_mitm.c} $\to$ \texttt{agents/AUD3/out/aud3\_mitm\_L25\_a8.txt}, 49\,s, 비트집합 $2.84$\,GB)
$\Lam=\Lam_{25}=26\,771\,144\,400$, $a_8=26\,577\,893\,791$이라 하면 $a_8$은 $\Lam$의 진약수 $7$개 이하(반복 허용)의 합이 아니다. 따라서 $\kappa_{25}(a_8)\ge8$이고 $H(25)\ge8$이다.
\end{proposition}
\begin{proof}[논법]
$a_8=d_1+\dots+d_k$ ($k\le7$, $d_1\ge\dots\ge d_k$)이면 $d_1\ge a_8/7$. $d_1=\Lam/m$이고 $m\le7\Lam/a_8=7.0509$이므로 $m\in\{2,\dots,7\}$ ($m=1$은 $\Lam>a_8$).
나머지는 $6$개 이하 약수의 합, 즉 $a_8-\Lam/m\in S_3+S_3$ ($S_3$: $3$개 이하 진약수의 합). 새 코드로 $S_3\cap[0,a_8-\Lam/7]$ ($|\,\cdot\,|=244\,906\,542$)을 만들고
여섯 수 $a_8-\Lam/m$ ($=13\,192\,321\,591$, $17\,654\,178\,991$, $19\,885\,107\,691$, $21\,223\,664\,911$, $22\,116\,036\,391$, $22\,753\,444\,591$) 각각에 대해
모든 $s\in S_3$, $s\le x$에서 $x-s\notin S_3$임을 전수 확인하였다(대조: $22\,753\,444\,590\in S_3+S_3$).
반복 허용 표현이 없으면 서로 다른 약수의 표현도 없으므로 중복 제거 정리조차 필요 없다.
\end{proof}
같은 방법으로 $L=23$에서 $\kappa(a_7)\ge7$ (BFS 독립, NUM의 \texttt{mitm23.c}와도 독립인 새 코드) 및 $\kappa(a_c)\ge c$ ($c=4,5,6$, $S_3+S_{c-4}$ 판정)와 $a_c-1\in R_{c-1}$을 확인하였다.
$H(23)\le7$, $H(25)\le8$ (모든 $a$가 덮임)은 여전히 BFS 두 구현의 일치에만 근거한다(NUM이 이미 밝힌 대로).

\paragraph{알고리즘 정당성.}
명제 \texttt{NUM:prop-bfs}는 (i) 반복 허용 합의 부분합이 모두 $[0,\Lam)$ 안에 있다는 것(절단 무손실)과 (ii) 중복 제거 정리만 쓴다. (ii)를 위에서 다시 유도하였다.
비트 이동 구현(\texttt{hbfs.c}, \texttt{hbfs2.c})은 감사자가 코드 수준으로 검증하지 않았으나, 결과는 AUD3의 독립 DP($L\le18$)와 독립 MITM($L=23,25$의 $R_3$ 크기와 표본 판정)으로 교차 확인된다.

\paragraph{NUM 수정 목록 (파일 내 \texttt{\% AUD3:}).}
(1) \texttt{prop-valid}(4)에 AUD3 독립 재검증 문장 추가. (2) \texttt{thm-H25}(4),(8)과 원장에 $\kappa(a_8)\ge8$의 BFS 독립 확인 추가.
(3) \texttt{tab-res} 설명에 LLL 점의 재현성 단서(유효숫자 15자리 저장) 추가. \texttt{frag\_src.tex}와 생성 스크립트는 고치지 않았다.

\paragraph{두 절 사이의 정합성.}
VKGAP와 NUM은 같은 두 수치를 다르게 적고 있었다: 반박 문턱(VKGAP $0.0194$, NUM $0.019428$; 정답 $1-|F(1)|=0.0194276$)과 극한 분산(VKGAP $0.03907$, NUM $0.039094$; 정답 $0.0390940$).
두 경우 모두 NUM 값이 옳고 VKGAP를 고쳤다. NUM 원장의 ``$\eta>0.0195$''는 (충분조건으로) 옳다. 두 절 모두 Erdős \#304와 SP1을 \OPEN{}으로 둔다.
